\chapter*{B1\quad 试题库（四）参考答案}
\addcontentsline{toc}{chapter}{B1\quad 试题库（四）参考答案}
\markboth{B1\quad 试题库（四）参考答案}{}

\begin{tishi}
B1 原卷的答案字母直接印在题干括号内，此处逐题列出并补一句判定理由（选择题），
问答题给出可作简答答案的完整要点，计算题按\textbf{已知 $\to$ 公式 $\to$ 代入 $\to$ 结果（带单位）}写全步骤。
答案册不标注星级，星级与依据见题目册每题末尾。
\end{tishi}

\section*{一、选择题（每题 2 分，共 30 分）}

\noindent\textbf{1. C\quad} 单位质量力＝质量力除以质量，即 $\bF/m$，量纲为加速度；
A 是单位面积力（压强），B 是单位体积力（重度），D 是单位重量流体受到的重力（无量纲的 1）。

\noindent\textbf{2. A\quad} 汽油的切应力与速度梯度成正比（$\tau=\mu\dfrac{\dd u}{\dd y}$，比例系数 $\mu$ 为常数），
是牛顿流体；纸浆、血液、沥青属非牛顿流体。

\noindent\textbf{3. C\quad} 由欧拉平衡微分方程 $\dd p=\rho(f_x\dd x+f_y\dd y+f_z\dd z)$ 可积，
要求体积力有势；其余三项都不是平衡的必要条件。

\noindent\textbf{4. D\quad} 欧拉法（空间点法）研究\textbf{通过空间各固定点的流体质点}的运动参数
随时间的变化；A、B 是拉格朗日法（质点法、迹线法）。

\noindent\textbf{5. C\quad} 定常（恒定）流中流线与迹线重合；
无旋、有旋只影响涡量，非定常流中流线与迹线一般不重合。

\noindent\textbf{6. C\quad} 沿总流的伯努利方程要求两个计算断面为\textbf{渐变流断面}，
但两断面\textbf{之间}允许出现急变流（其间的水头损失另行计入）；A、B、D 都说得绝对。

\noindent\textbf{7. B\quad} 容器自由下落（完全失重）时质量力只剩惯性力且处处为零，
由 $\dd p=\rho(f_x\dd x+f_y\dd y+f_z\dd z)=0$ 得 $p$ 处处相等，故 $h$ 处压强仍为 $p_0$。

\noindent\textbf{8. C\quad} 虹吸管最高处要靠管内负压才能把水"提"上去，
故最高处压强小于大气压（这也是限制虹吸安装高度的原因）。

\noindent\textbf{9. A\quad} 总水头线沿程只降不升（沿程必有水头损失）；
测压管水头线可升可降（管径变化使流速水头与压强水头互相转换），故 A 正确。

\noindent\textbf{10. B\quad} 两管断面面积、长度、相对粗糙度相同且都处于紊流粗糙区，
$\lambda$ 只取决于相对粗糙度；由 $h_f=\lambda\dfrac{l}{4R}\dfrac{v^2}{2g}$ 且 $v=Q/A$，
水力半径之比 $R_{\text{圆}}:R_{\text{方}}=1.128:1$，
故 $\dfrac{h_{f\text{圆}}}{h_{f\text{方}}}\approx0.886$。

\noindent\textbf{11. C\quad} 运动黏滞系数 $\nu=\mu/\rho$，量纲为 $\mathrm{L^2/T}$；
A 为加速度、B 为 $\mathrm{L/T^3}$、D 为速度量纲。

\noindent\textbf{12. B\quad} 有压管流中黏滞阻力起主导作用，模型试验应满足\textbf{雷诺准则}；
A 弗劳德准则用于明渠等重力主导的流动。

\noindent\textbf{13. D\quad} 无旋运动的判据是涡量（旋转角速度的 2 倍）为零；
角变形为零只是无角变形，与无旋不是一回事；流线、迹线是否为直线也与无旋无关。

\noindent\textbf{14. A\quad}（流体机械）允许吸上真空高度 $H_s$ 表征泵入口处允许达到的最大真空度，
反映的是离心泵的\textbf{吸水性能}。

\noindent\textbf{15. C\quad}（流体机械）两台同型号风机并联，并联风量小于单台风量之和，
即\textbf{大幅度增加但不是成倍增加}。

\section*{二、问答题（每题 6 分，共 12 分）}

\noindent\textbf{第 1 题\quad 流动性与黏滞性、连续介质模型}

\begin{jie}
\textbf{1.\ 流动性}

流动性是指流体\textbf{在微小的切力作用下就会产生连续变形（流动），且只要切力存在变形就持续进行}，
即流体不能保持固定的形状、只能保持一定的体积。流体分子间内聚力小、分子易相对滑移，
是流动性的根本原因；流动性使流体只能承受压力而几乎不能承受拉力与切力。

\textbf{2.\ 黏滞性}

黏滞性是指流体\textbf{内部相邻流层之间因相对运动而产生内摩擦阻力（切应力）}的性质，
其大小由牛顿内摩擦定律 $\tau=\mu\dfrac{\dd u}{\dd y}$ 描述，$\mu$ 为动力黏度。
黏滞性使流体在运动中产生机械能损失（水头损失），是实际流体区别于理想流体的根本标志。

\textbf{3.\ 两者的关系}

流动性使流体不断变形（"会流"），黏滞性使变形过程中产生阻力（"不好流"），
二者是流体同一对矛盾的两种表现：理想流体保留了流动性而去掉了黏滞性。

\textbf{4.\ 连续介质模型}

把流体视为由\textbf{无数连续分布的流体质点}组成的介质，其物理量（密度 $\rho$、速度 $u$、
压强 $p$ 等）是空间坐标与时间的\textbf{连续函数}。其中"质点"指宏观上足够小（可视为几何点，
内部物理量均匀）、微观上足够大（含大量分子，统计平均有确定值）的流体微团。

引入该模型的必要性：① 微积分、微分方程要求所研究的量连续可导，而真实流体由离散分子组成，
直接按分子描述无法求导、取极限；② 取统计平均后可得确定的宏观量；
③ 工程中特征尺度远大于分子平均自由程，抹平分子间隙带来的误差可忽略。

\textbf{适用条件}：所研究流动的特征尺度远大于分子平均自由程；高空稀薄气体、真空技术不适用。
\end{jie}

\noindent\textbf{第 2 题\quad 泵与风机的标准条件与相似换算（流体机械，复试范围）}

\begin{jie}
\textbf{1.\ 风机的标准条件}

厂家样本参数是在标准进气状态下试验得到的，一般风机的标准条件为：
\textbf{大气压 $p_0=101.325\ \mathrm{kPa}$（1 个标准大气压）}、
\textbf{温度 $t_0=20\ ^{\circ}\mathrm{C}$}、
\textbf{相对湿度 $\varphi=50\%$}、
\textbf{气体密度 $\rho_0=1.2\ \mathrm{kg/m^3}$}。

\textbf{2.\ 相似换算关系式}

当实际使用条件与样本条件不同时，按相似律（同一台泵或风机、转速不变）换算。
因流量与效率只取决于几何与转速，与气体状态无关：

\[ Q=Q_0,\qquad \eta=\eta_0 \]

压力（全压或静压）与气体密度成正比，而 $\rho\propto\dfrac{B}{T}=\dfrac{B}{273+t}$，故

\[ p=p_0\frac{B}{101325}\cdot\frac{293}{273+t} \]

轴功率 $N=\dfrac{pQ}{\eta}$，由 $p$ 的换算式得

\[ N=N_0\frac{B}{101325}\cdot\frac{293}{273+t} \]

即：\textbf{流量与效率不随气体状态改变}；\textbf{全压与轴功率均正比于气体密度}，
即
\[ \frac{N}{N_0}=\frac{p}{p_0}=\frac{\rho}{\rho_0} \]

\textbf{说明}：若同时改变转速，还要叠加转速相似律
$\dfrac{Q}{Q_0}=\dfrac{n}{n_0}$、$\dfrac{p}{p_0}=\left(\dfrac{n}{n_0}\right)^2$、
$\dfrac{N}{N_0}=\left(\dfrac{n}{n_0}\right)^3$。本题按"转速不变、只改变进气状态"作答。
\end{jie}

\section*{三、计算题（共 58 分）}

\noindent\textbf{第 1 题\quad 弧形闸门静水总压力与力矩（15 分）}

\begin{jie}
\textbf{已知}：闸门宽 $B=2\ \mathrm{m}$，圆弧半径 $r=3\ \mathrm{m}$，圆心角 $\alpha=30^{\circ}$；
转轴 $O$ 与水面齐平，且为圆弧的圆心（题给"闸门转轴与水平面齐平"即指 $O$ 位于水面）；
水的重度 $\gamma=\rho g=9.8\times10^{3}\ \mathrm{N/m^3}$。

\textbf{1.\ 几何关系}

取 $O$ 为原点、水面为基准面；弧面以圆心角 $\theta$ 参数化，$\theta$ 自水平方向（$A$ 点所在半径）
起算、取 $0\sim\alpha=30^{\circ}$，$z$ 轴向下为正。弧上任一点：

\[ x=r\cos\theta,\qquad z=r\sin\theta,\qquad p=\gamma z=\gamma r\sin\theta \]

弧的竖直投影高度与水平投影长度：

\[ h_z=r\sin\alpha=3\times0.5=1.5\ \mathrm{m},\qquad
h_x=r(1-\cos\alpha)=3\times(1-0.8660)=0.402\ \mathrm{m} \]

\textbf{2.\ 水平分力 $F_x$}

方法一（投影面法）：曲面上静水总压力的水平分力等于曲面在铅垂面上投影所得的矩形平面
所受的静水总压力，该矩形高 $r\sin\alpha$、宽 $B$，形心淹没深度为 $\dfrac{r\sin\alpha}{2}$：

\[ F_x=\gamma\times B\,r\sin\alpha\times\frac{r\sin\alpha}{2}
=\gamma B\frac{(r\sin\alpha)^2}{2} \]

代入：

\[ F_x=9.8\times10^{3}\times2\times\frac{(3\times0.5)^2}{2}
=9.8\times10^{3}\times2.25=2.205\times10^{4}\ \mathrm{N}=22.05\ \mathrm{kN} \]

方法二（沿弧面积分）：弧微元长度 $r\dd\theta$，其压力在水平方向的分量为
$p\,r\dd\theta\cdot\cos\theta$：

\[ F_x=\int_0^{\alpha}\gamma r\sin\theta\cdot B\,r\dd\theta\cdot\cos\theta
=\gamma Br^2\int_0^{\alpha}\sin\theta\cos\theta\,\dd\theta
=\gamma Br^2\left[\frac{\sin^2\theta}{2}\right]_0^{\alpha}
=\frac{\gamma Br^2}{2}\sin^2\alpha \]

\[ F_x=\frac{9.8\times10^{3}\times2\times9}{2}\times0.25=9.8\times10^{3}\times2.25
=2.205\times10^{4}\ \mathrm{N}=22.05\ \mathrm{kN} \]

两种方法结果一致。方向：\textbf{水平指向闸门内侧}（由水体指向闸门）。

\textbf{3.\ 竖直分力 $F_z$}

方法一（压力体法）：竖直分力等于曲面上方到水面之间"压力体"内液体的重力。
压力体由弧面 $AB$、半径 $OA$ 与 $OB$ 围成，面积为扇形 $OAB$ 减去三角形 $OAB$：

\[ A_{\text{扇}}=\frac{\alpha}{2}r^2=\frac{\pi/6}{2}\times9=2.356\ \mathrm{m^2},\qquad
A_{\triangle}=\frac{1}{2}r^2\sin\alpha=\frac{1}{2}\times9\times0.5=2.25\ \mathrm{m^2} \]

\[ A_{\text{压}}=2.356-2.25=0.106\ \mathrm{m^2},\qquad
F_z=\gamma B A_{\text{压}}=9.8\times10^{3}\times2\times0.106
=2.08\times10^{3}\ \mathrm{N}=2.08\ \mathrm{kN} \]

方法二（沿弧面积分）：压力在竖直方向的分量为 $p\,r\dd\theta\cdot\sin\theta$：

\[ F_z=\int_0^{\alpha}\gamma r\sin\theta\cdot B\,r\dd\theta\cdot\sin\theta
=\gamma Br^2\int_0^{\alpha}\sin^2\theta\,\dd\theta
=\gamma Br^2\left[\frac{\theta}{2}-\frac{\sin2\theta}{4}\right]_0^{\alpha} \]

\[ \int_0^{\pi/6}\sin^2\theta\,\dd\theta=\frac{\pi}{12}-\frac{\sqrt3}{8}
=0.2618-0.21651=0.0453 \]

\[ F_z=9.8\times10^{3}\times3^{2}\times2\times0.0453
=9.8\times10^{3}\times2\times0.4077
=7.99\times10^{3}\ \mathrm{N}\approx8.0\ \mathrm{kN} \]

两种方法的结果不同（$2.08$ 与 $8.0\ \mathrm{kN}$），原因是压力体的边界取法：
按本图几何，弧面 $AB$ 下方为水体、弧面两侧均与水面连通，
真正的压力体应是\textbf{弧面以上的水体}，即由弧面、水面（经 $A$ 点）与 $OB$ 围成的区域，
其面积为

\[ A_{\text{压}}=\frac{r^2}{2}\left(\alpha-\sin\alpha\cos\alpha\right)
=\frac{9}{2}\left(0.5236-0.5\times0.8660\right)=4.5\times0.0906=0.4077\ \mathrm{m^2} \]

\[ F_z=\gamma B A_{\text{压}}=9.8\times10^{3}\times2\times0.4077
=7.99\times10^{3}\ \mathrm{N}=7.99\ \mathrm{kN} \]

故以\textbf{积分（压力体几何）结果}为准：

\[ F_z=7.99\ \mathrm{kN}\ (\text{竖直向下}) \]

\textbf{4.\ 对转轴 $O$ 的总合力矩 $M$}

圆柱面上任一点的静水压强沿\textbf{曲面法线（半径方向）}作用，其作用线\textbf{全部通过圆心 $O$}，
对 $O$ 点的力臂恒为零，故

\[ M=\int_{\text{弧面}}\dd F\times0=0 \]

即作用于弧形闸门上的静水总压力对转轴（圆心）的总力矩为 $M=0$。

\textbf{5.\ 结果}

\[ F_x=22.05\ \mathrm{kN}\ (\text{水平，指向闸门内侧}),\qquad
F_z=7.99\ \mathrm{kN}\ (\text{竖直向下}),\qquad M=0 \]

\begin{yicuo}
\textbf{易错点：}① 只有 $O$ 是圆弧圆心时，压力作用线才全部过 $O$、力矩才为零；
② 水平分力是"投影面形心压强 $\times$ 投影面积"，竖直投影高度是 $r\sin\alpha=1.5\ \mathrm{m}$
而不是 $r$；③ 竖直分力的压力体是\textbf{弧面以上直到水面的那块水体}，
不是扇形与三角形之差，也不是整个半圆。
\end{yicuo}
\end{jie}

\noindent\textbf{第 2 题\quad 水箱垂直管道出流的流量与管长关系（15 分）}

\begin{jie}
\textbf{已知}：水箱水深（作用水头）$H$，管径 $d$，管长 $l$，沿程阻力系数 $\lambda$，
局部阻力系数 $\zeta$；管道向大气出流。

\textbf{1.\ 列能量方程}

取水箱水面为断面 1-1、管道出口为断面 2-2，基准面取出口所在水平面。
水面与出口均通大气，$p_1=p_2=p_a$；水面面积远大于管道，$v_1\approx0$；
出口流速为 $v$。水面到出口的高程差即为水箱水深 $H$，故 $z_1=H$、$z_2=0$。
水头损失由进口局部损失与沿程损失组成：

\[ h_w=\left(\zeta+\lambda\frac{l}{d}\right)\frac{v^2}{2g} \]

能量方程：

\[ H+0+0=0+0+\frac{v^2}{2g}+\left(\zeta+\lambda\frac{l}{d}\right)\frac{v^2}{2g} \]

整理：

\[ \frac{v^2}{2g}=\frac{H}{1+\zeta+\lambda\dfrac{l}{d}} \]

\textbf{2.\ 流量表达式}

\[ v=\sqrt{\frac{2gH}{1+\zeta+\lambda l/d}},\qquad
Q=\frac{\pi d^2}{4}v=\frac{\pi d^2}{4}\sqrt{\frac{2gH}{1+\zeta+\lambda l/d}} \]

\textbf{3.\ 对 $l$ 的变化规律}

分子 $2gH$ 与 $l$ 无关，分母 $1+\zeta+\lambda\dfrac{l}{d}$ 随 $l$ 单调增大，故

\[ \frac{\dd Q}{\dd l}<0 \]

即\textbf{管长越大、流量越小}。回到题目所问的三种情形：

\begin{xiaowen}
\item $Q$ 不随管长 $l$ 而变的条件：要求分母与 $l$ 无关，即
$\dfrac{\lambda}{d}=0$，也就是 $\lambda=0$（\textbf{无沿程损失，即理想流体}）。
工程上只要 $\lambda l/d\ll1+\zeta$（管很短或沿程阻力很小），$Q$ 就近似与 $l$ 无关；
\textbf{只有 $\lambda=0$ 时严格不变}。
\item $Q$ 随 $l$ 加大而增加的条件：由上式\textbf{不可能出现}；
除非水位随管长增加而提升（$H$ 随 $l$ 增大），此时
$\dfrac{\dd Q}{\dd l}>0$ 要求 $\dfrac{\dd H}{\dd l}>H\dfrac{\lambda/d}{1+\zeta+\lambda l/d}$。
对本问题给定的固定 $H$，这一情形不存在。
\item $Q$ 随 $l$ 加大而减小的条件：只要 $\lambda>0$（实际流体），
\textbf{恒成立}。由式还可得定量关系：
$\dfrac{\dd Q/\dd l}{Q}=-\dfrac{\lambda}{2d\left(1+\zeta+\lambda l/d\right)}$。
\end{xiaowen}

\textbf{4.\ 讨论：原题若要三问都成立}

原题把三问并列，说明其模型的"驱动水头"本身含 $l$。
若把管道理解为\textbf{自水面竖直向下}、出口在管末、水面到出口的高程差为 $H+l$
（即水位高出管道进口 $H$、管道再向下延伸 $l$），则

\[ \frac{v^2}{2g}=\frac{H+l}{1+\zeta+\lambda\dfrac{l}{d}},\qquad
Q=\frac{\pi d^2}{4}\sqrt{\frac{2g(H+l)}{1+\zeta+\lambda l/d}} \]

记 $f(l)=\dfrac{H+l}{1+\zeta+\lambda l/d}$，求导：

\[ f'(l)=\frac{\left(1+\zeta+\lambda l/d\right)-\dfrac{\lambda}{d}(H+l)}
{\left(1+\zeta+\lambda l/d\right)^2}
=\frac{1+\zeta-\dfrac{\lambda H}{d}}{\left(1+\zeta+\lambda l/d\right)^2} \]

\textbf{$f'(l)$ 的符号与 $l$ 无关}，只取决于 $1+\zeta-\dfrac{\lambda H}{d}$ 的符号，
于是三问恰好对应：

\[ \begin{cases}
1+\zeta>\dfrac{\lambda H}{d}, & f'(l)>0,\ Q\ \text{随}\ l\ \text{加大而增加}\\[6pt]
1+\zeta=\dfrac{\lambda H}{d}, & f'(l)=0,\ Q\ \text{与}\ l\ \text{无关（不变）}\\[6pt]
1+\zeta<\dfrac{\lambda H}{d}, & f'(l)<0,\ Q\ \text{随}\ l\ \text{加大而减小}
\end{cases} \]

\textbf{5.\ 结论}

按本册采用的基准（水面到出口高程差为水深 $H$）：

\[ Q=\frac{\pi d^2}{4}\sqrt{\frac{2gH}{1+\zeta+\lambda l/d}} \]

$Q$ 严格不随 $l$ 变的唯一条件是 $\lambda=0$；$Q$ 随 $l$ 增大只减不增；
而原题"$Q$ 随 $l$ 加大而增加"一问，只有在把"$H+l$"作为驱动水头、
且满足 $1+\zeta>\lambda H/d$ 时才成立。考生作答时把两种基准都写出并说明所选基准即可。
\end{jie}

\noindent\textbf{第 3 题\quad 煤粉颗粒的沉降与被带走的判别（12 分）}

\begin{jie}
\textbf{已知}：烟气上升速度 $v=0.5\ \mathrm{m/s}$；
烟气运动黏滞系数 $\nu=230\times10^{-6}\ \mathrm{m^2/s}$；烟气密度 $\rho=0.2\ \mathrm{kg/m^3}$；
煤粉密度 $\rho_m=1.3\times10^{3}\ \mathrm{kg/m^3}$；
颗粒直径 $d=0.1\ \mathrm{mm}=1\times10^{-4}\ \mathrm{m}$。

\textbf{1.\ 先求烟气的动力黏度}

题给的是运动黏度，沉降公式中要用动力黏度：

\[ \mu=\rho\nu=0.2\times230\times10^{-6}=4.6\times10^{-5}\ \mathrm{Pa\cdot s} \]

\textbf{2.\ 按斯托克斯（层流）区求沉降速度}

\[ v_t=\frac{d^2(\rho_m-\rho)g}{18\mu} \]

代入：

\[ v_t=\frac{(1\times10^{-4})^2\times(1.3\times10^{3}-0.2)\times9.8}
{18\times4.6\times10^{-5}}
=\frac{1\times10^{-8}\times1299.8\times9.8}{8.28\times10^{-4}} \]

\[ v_t=\frac{1.274\times10^{-4}}{8.28\times10^{-4}}=0.154\ \mathrm{m/s} \]

\textbf{3.\ 校核雷诺数}

\[ \Rey_p=\frac{v_t d}{\nu}=\frac{0.154\times1\times10^{-4}}{230\times10^{-6}}
=0.067<1 \]

满足斯托克斯区的适用条件（$\Rey_p<1$），故上述沉降速度成立，
不必改用过渡区或牛顿阻力区公式。

\textbf{4.\ 与烟气流速比较，判定沉降还是被带走}

颗粒在炉膛中的绝对速度（向上为正）＝烟气上升速度 $-$ 颗粒沉降速度：

\[ v_{\text{净}}=v-v_t=0.5-0.154=0.346\ \mathrm{m/s}>0 \]

即颗粒以 $0.346\ \mathrm{m/s}$ 的净速度\textbf{随烟气向上运动}，
故煤粉颗粒\textbf{被烟气带走}，不会沉降下来。

\textbf{5.\ 结论}

\[ v_t\approx0.154\ \mathrm{m/s}<v=0.5\ \mathrm{m/s}
\ \Longrightarrow\ \text{颗粒被上升烟气带走} \]

\begin{yicuo}
\textbf{易错点：}① 公式中的黏度必须是\textbf{动力黏度} $\mu$，题给的是运动黏度 $\nu$，
要先换算 $\mu=\rho\nu$；若误用 $\nu$ 代入会被放大到约 $0.15$ 与 $0.03$ 相差一个数量级；
② 算出沉降速度后必须用 $\Rey_p=\dfrac{v_td}{\nu}$ 校核是否在层流区；
③ 判据是"沉降速度与烟气流速比大小"，不是"颗粒密度比气体大就会沉降"。
\end{yicuo}
\end{jie}

\noindent\textbf{第 4 题\quad 水泵工况点与节流调节（16 分，流体机械）}

\begin{jie}
\textbf{已知}：$n=1450\ \mathrm{r/min}$；性能参数表（$Q$ 单位 $\mathrm{L/s}$）：

\begin{center}
\begin{tabular}{@{}lccccccccc@{}}
\toprule
$Q$ & 0 & 2 & 4 & 6 & 8 & 10 & 12 & 14 \\
\midrule
$H$（$\mathrm{m}$） & 11 & 10.8 & 10.5 & 10 & 9.2 & 8.4 & 7.4 & 6 \\
$\eta$（\%） & 0 & 15 & 30 & 45 & 60 & 65 & 55 & 30 \\
\bottomrule
\end{tabular}
\end{center}

管路综合阻力数 $S=0.024\times10^{3}\ \mathrm{s^2/m^5}$；几何扬水高度 $H_Z=6\ \mathrm{m}$；
上下水池水面均为大气压；水的密度 $\rho=1000\ \mathrm{kg/m^3}$。

\textbf{1.\ 管路特性曲线}

\[ h=H_Z+SQ^2=6+0.024\times10^{3}Q^2 \qquad (Q\ \text{以}\ \mathrm{m^3/s}\ \text{计}) \]

把 $Q$ 换成以 $\mathrm{L/s}$ 计（$Q\ [\mathrm{m^3/s}]=10^{-3}Q\ [\mathrm{L/s}]$）：

\[ h=6+0.024\times10^{3}\times10^{-6}Q^2=6+0.024Q^2 \qquad (Q\ \text{以}\ \mathrm{L/s}\ \text{计}) \]

\textbf{2.\ 求工况点（泵性能曲线与管路特性曲线的交点）}

逐点比较：

\begin{center}
\begin{tabular}{@{}lcccccccc@{}}
\toprule
$Q$（$\mathrm{L/s}$） & 0 & 2 & 4 & 6 & 8 & 10 & 12 & 14 \\
\midrule
管路需要水头 $h$（$\mathrm{m}$） & 6 & 6.10 & 6.38 & 6.86 & 7.54 & 8.40 & 9.46 & 10.70 \\
泵提供水头 $H$（$\mathrm{m}$） & 11 & 10.8 & 10.5 & 10 & 9.2 & 8.4 & 7.4 & 6 \\
\midrule
差值 $H-h$（$\mathrm{m}$） & 5 & 4.70 & 4.12 & 3.14 & 1.66 & 0 & $-2.06$ & $-4.70$ \\
\bottomrule
\end{tabular}
\end{center}

差值由 $Q=8\ \mathrm{L/s}$ 处的 $+1.66\ \mathrm{m}$ 变为 $Q=10\ \mathrm{L/s}$ 处的 $0$，
故交点在 $Q=8\sim10\ \mathrm{L/s}$ 之间。在 $Q=9.4\ \mathrm{L/s}$ 处插值：
$H\approx9.2-\dfrac{1.4}{2}\times0.8=8.64\ \mathrm{m}$，$h=6+0.024\times9.4^2=8.12\ \mathrm{m}$（还差 $0.5$）；
在 $Q=9.6$ 处 $H\approx8.56$、$h=8.21$（差 $0.35$）。继续逼近得工况点

\[ Q\approx9.9\ \mathrm{L/s},\qquad H=6+0.024\times9.9^2\approx8.35\ \mathrm{m} \]

取整数化结果

\[ Q\approx10\ \mathrm{L/s},\qquad H=8.4\ \mathrm{m},\qquad \eta=65\% \]

（说明：若按上表逐点比较，$Q=10\ \mathrm{L/s}$ 处泵提供水头 $H=8.4\ \mathrm{m}$
恰等于管路需要水头 $h=6+0.024\times100=8.4\ \mathrm{m}$，
\textbf{该点正好是两曲线的交点}，故工况点就是表格中的 $Q=10\ \mathrm{L/s}$ 这一列。）

\textbf{3.\ 求轴功率}

\[ N=\frac{\rho g QH}{\eta}
=\frac{1000\times9.8\times10\times10^{-3}\times8.4}{0.65}
=\frac{823.2}{0.65}=1.266\times10^{3}\ \mathrm{W} \]

\[ Q=10\ \mathrm{L/s},\quad H=8.4\ \mathrm{m},\quad \eta=65\%,\quad
N\approx1.27\times10^{3}\ \mathrm{W}=1.27\ \mathrm{kW} \]

\textbf{4.\ 运行工况的调节方式}

\begin{xiaowen}
\item \textbf{节流调节}：关小出口阀门，增大管路阻力（$S$ 增大），工况点沿泵性能曲线
向小流量方向移动；设备简单，但额外消耗水头，不经济；
\item \textbf{变速调节}：改变水泵转速 $n$，按相似律
$\dfrac{Q}{Q_0}=\dfrac{n}{n_0}$、$\dfrac{H}{H_0}=\left(\dfrac{n}{n_0}\right)^2$、
$\dfrac{N}{N_0}=\left(\dfrac{n}{n_0}\right)^3$ 改变性能曲线，最节能，应用最广；
\item \textbf{切削叶轮（换轮）调节}：改变叶轮外径，性能曲线按切削律平移；
\item \textbf{管路系统调节}：改变管径、管长或开启并联支路（改变 $S$）；
\item \textbf{台数调节}：多台泵并联或串联运行，改变系统总性能曲线。
\end{xiaowen}

\textbf{5.\ 节流调节到 $Q=8\ \mathrm{L/s}$ 时的参数}

由性能表，$Q=8\ \mathrm{L/s}$ 时泵的

\[ H=9.2\ \mathrm{m},\qquad \eta=60\% \]

此时管路系统所需水头 $h=6+0.024\times8^2=6+1.536=7.54\ \mathrm{m}$，
多余水头由节流阀消耗：

\[ h_{\text{阀}}=9.2-7.54=1.66\ \mathrm{m} \]

轴功率：

\[ N=\frac{\rho gQH}{\eta}
=\frac{1000\times9.8\times8\times10^{-3}\times9.2}{0.60}
=\frac{721.3}{0.60}=1.202\times10^{3}\ \mathrm{W} \]

\[ Q=8\ \mathrm{L/s},\quad H=9.2\ \mathrm{m},\quad \eta=60\%,\quad
N\approx1.20\times10^{3}\ \mathrm{W}=1.20\ \mathrm{kW} \]

\textbf{6.\ 结果汇总与讨论}

\[ \text{工况点：}\ Q=10\ \mathrm{L/s},\ H=8.4\ \mathrm{m},\ \eta=65\%,\
N\approx1.27\times10^{3}\ \mathrm{W} \]
\[ \text{节流至}\ Q=8\ \mathrm{L/s}:\ H=9.2\ \mathrm{m},\ \eta=60\%,\
N\approx1.20\times10^{3}\ \mathrm{W} \]

流量下降 $20\%$、轴功率下降约 $5\%$，另有多余水头 $1.66\ \mathrm{m}$ 被阀门白白消耗，
故选 \textbf{变速调节}优于节流调节。
\end{jie}
